\documentclass[10pt,a4paper]{amsart}
\usepackage{amsmath,amssymb,amsthm,graphicx}
% Retained minimal preamble for the Stacks Project chapter more-morphisms.tex (More on Morphisms).
% Verbatim excerpt of the chapter preamble: Xy-pic setup (native preamble lines 14-18) and the native
% theorem declarations with the mathematical macros (native preamble lines 154-217). Copied mechanically;
% nothing added, removed or redefined. The wrapper class and the amsmath/amssymb/amsthm/graphicx base
% packages are supplied by the assembler, so the class/externaldocument part of the native preamble is omitted.
% For commutative diagrams we use Xy-pic
\usepackage[all]{xy}

% We use 2cell for 2-commutative diagrams.
\xyoption{2cell}
% Theorem environments.
%
\theoremstyle{plain}
\newtheorem{theorem}[subsection]{Theorem}
\newtheorem{proposition}[subsection]{Proposition}
\newtheorem{lemma}[subsection]{Lemma}

\theoremstyle{definition}
\newtheorem{definition}[subsection]{Definition}
\newtheorem{example}[subsection]{Example}
\newtheorem{exercise}[subsection]{Exercise}
\newtheorem{situation}[subsection]{Situation}

\theoremstyle{remark}
\newtheorem{remark}[subsection]{Remark}
\newtheorem{remarks}[subsection]{Remarks}

\numberwithin{equation}{subsection}

% Macros
%
\def\lim{\mathop{\mathrm{lim}}\nolimits}
\def\colim{\mathop{\mathrm{colim}}\nolimits}
\def\Spec{\mathop{\mathrm{Spec}}}
\def\Hom{\mathop{\mathrm{Hom}}\nolimits}
\def\Ext{\mathop{\mathrm{Ext}}\nolimits}
\def\SheafHom{\mathop{\mathcal{H}\!\mathit{om}}\nolimits}
\def\SheafExt{\mathop{\mathcal{E}\!\mathit{xt}}\nolimits}
\def\Sch{\mathit{Sch}}
\def\Mor{\mathop{\mathrm{Mor}}\nolimits}
\def\Ob{\mathop{\mathrm{Ob}}\nolimits}
\def\Sh{\mathop{\mathit{Sh}}\nolimits}
\def\NL{\mathop{N\!L}\nolimits}
\def\CH{\mathop{\mathrm{CH}}\nolimits}
\def\proetale{{pro\text{-}\acute{e}tale}}
\def\etale{{\acute{e}tale}}
\def\QCoh{\mathit{QCoh}}
\def\Ker{\mathop{\mathrm{Ker}}}
\def\Im{\mathop{\mathrm{Im}}}
\def\Coker{\mathop{\mathrm{Coker}}}
\def\Coim{\mathop{\mathrm{Coim}}}

% Boxtimes
%
\DeclareMathSymbol{\boxtimes}{\mathbin}{AMSa}{"02}

%
% Macros for moduli stacks/spaces
%
\def\QCohstack{\mathcal{QC}\!\mathit{oh}}
\def\Cohstack{\mathcal{C}\!\mathit{oh}}
\def\Spacesstack{\mathcal{S}\!\mathit{paces}}
\def\Quotfunctor{\mathrm{Quot}}
\def\Hilbfunctor{\mathrm{Hilb}}
\def\Curvesstack{\mathcal{C}\!\mathit{urves}}
\def\Polarizedstack{\mathcal{P}\!\mathit{olarized}}
\def\Complexesstack{\mathcal{C}\!\mathit{omplexes}}
% \Pic is the operator that assigns to X its picard group, usage \Pic(X)
% \Picardstack_{X/B} denotes the Picard stack of X over B
% \Picardfunctor_{X/B} denotes the Picard functor of X over B
\def\Pic{\mathop{\mathrm{Pic}}\nolimits}
\def\Picardstack{\mathcal{P}\!\mathit{ic}}
\def\Picardfunctor{\mathrm{Pic}}
\def\Deformationcategory{\mathcal{D}\!\mathit{ef}}
\newenvironment{reference}{}{}
\newenvironment{slogan}{}{}
\newenvironment{history}{}{}
\usepackage{cleveref}
\crefname{originallemma}{original lemma}{original lemmas}
\Crefname{originallemma}{Original lemma}{Original lemmas}
\makeatletter
\newlabel{lemma-algebraically-independent}{{lemma-algebraically-independent}{source}{}{sourcepointer-1}{}}
\newlabel{lemma-algebraically-independent@cref}{{[originallemma][0][]lemma-algebraically-independent}{[1][0][]source}{}{}{}}
\makeatother
\crefname{originallemma}{original lemma}{original lemmas}
\Crefname{originallemma}{Original lemma}{Original lemmas}
\makeatletter
\newlabel{lemma-amazing-bertini-lemma}{{lemma-amazing-bertini-lemma}{source}{}{sourcepointer-2}{}}
\newlabel{lemma-amazing-bertini-lemma@cref}{{[originallemma][0][]lemma-amazing-bertini-lemma}{[1][0][]source}{}{}{}}
\makeatother
\crefname{originallemma}{original lemma}{original lemmas}
\Crefname{originallemma}{Original lemma}{Original lemmas}
\makeatletter
\newlabel{lemma-geom-irreducible-generic-fibre}{{lemma-geom-irreducible-generic-fibre}{source}{}{sourcepointer-3}{}}
\newlabel{lemma-geom-irreducible-generic-fibre@cref}{{[originallemma][0][]lemma-geom-irreducible-generic-fibre}{[1][0][]source}{}{}{}}
\makeatother
\crefname{originallemma}{original lemma}{original lemmas}
\Crefname{originallemma}{Original lemma}{Original lemmas}
\makeatletter
\newlabel{varieties-lemma-geometrically-irreducible-function-field}{{varieties-lemma-geometrically-irreducible-function-field}{source}{}{sourcepointer-4}{}}
\newlabel{varieties-lemma-geometrically-irreducible-function-field@cref}{{[originallemma][0][]varieties-lemma-geometrically-irreducible-function-field}{[1][0][]source}{}{}{}}
\makeatother
\crefname{originallemma}{original lemma}{original lemmas}
\Crefname{originallemma}{Original lemma}{Original lemmas}
\makeatletter
\newlabel{varieties-situation-family-divisors}{{varieties-situation-family-divisors}{source}{}{sourcepointer-5}{}}
\newlabel{varieties-situation-family-divisors@cref}{{[originallemma][0][]varieties-situation-family-divisors}{[1][0][]source}{}{}{}}
\makeatother
\title{A generic member of a sufficiently ample linear system stays geometrically irreducible}
\author{The Stacks Project Authors}
\date{stacks-project:more-morphisms.tex:local-snapshot-20261011; GNU FDL 1.2 or later}
\begin{document}
\maketitle
\noindent\textit{Editorial note.} Statement span is lines 9054-9071 of sources/stacks-more-morphisms/pinned-original-source.tex; proof span is lines 9073-9159, the complete author proof (85 body lines).\par\medskip
\noindent\textit{Editorial note.} Navigation stubs: 5 cited-but-not-included labels are supplied as non-load-bearing pointers with their pinned tags where the snapshot tag index has them.\par\medskip
\noindent\textit{Editorial note.} extra\_packages is empty on purpose: the retained preamble already loads xy with options and a second declaration clashes in the assembler wrapper.\par\medskip
\noindent\textit{Editorial note.} Text is retained verbatim, including the author's own typos and placeholders.\par\medskip
\noindent\textit{Editorial note.} The author's statement contains 'in reference \{\textbackslash{}cite[Theorem 6.3 part 4)]\{Jou\}\}', so the card carries a bibliography fragment. That fragment reproduces, field for field and without editing, the @BOOK\{Jou\} entry from the Stacks Project snapshot bibliography shipped with this corpus at focused-seeds/AGP/sources/stacks-snapshot-20261011/my.bib (same upstream Stacks Project snapshot). No entry text was altered and no citation was reconstructed from memory.\par\medskip
\noindent\textit{Editorial note.} A three-line wrapper preamble span supplies the mechanical reference/slogan/history environments, which the retained chapter preamble does not define; this follows the pattern already used by other Stacks cards in this corpus.\par\medskip
\noindent\textit{Allowed standard prerequisites (editorial).}\begin{itemize}
\item Explicit prerequisite boundary. Ample line bundles and their linear systems, geometric irreducibility, and generic flatness; the cited inputs are prerequisites and the generic-member argument is retained.
\end{itemize}\medskip
\ifdefined\hypertarget\hypertarget{sourcepointer-1}{}\fi
\noindent\textit{Source reference.} Original lemma lemma-algebraically-independent: sources/stacks-more-morphisms/pinned-original-source.tex, source lines 9012--9022. Cited by the retained statement or proof but not part of the collected spans; supplied only as a navigation pointer. The argument does not depend on its content.\par\medskip
\ifdefined\hypertarget\hypertarget{sourcepointer-2}{}\fi
\noindent\textit{Source reference.} Original lemma lemma-amazing-bertini-lemma: sources/stacks-more-morphisms/pinned-original-source.tex, source lines 8907--8910. Cited by the retained statement or proof but not part of the collected spans; supplied only as a navigation pointer. The argument does not depend on its content.\par\medskip
\ifdefined\hypertarget\hypertarget{sourcepointer-3}{}\fi
\noindent\textit{Source reference.} Original lemma lemma-geom-irreducible-generic-fibre: sources/stacks-more-morphisms/pinned-original-source.tex, source lines 7557--7564. Cited by the retained statement or proof but not part of the collected spans; supplied only as a navigation pointer. The argument does not depend on its content.\par\medskip
\ifdefined\hypertarget\hypertarget{sourcepointer-4}{}\fi
\noindent\textit{Source reference.} Original lemma varieties-lemma-geometrically-irreducible-function-field: sources/stacks-more-morphisms/pinned-original-source.tex, source lines 9159--9159. Cited by the retained statement or proof but not part of the collected spans; supplied only as a navigation pointer. The argument does not depend on its content.\par\medskip
\ifdefined\hypertarget\hypertarget{sourcepointer-5}{}\fi
\noindent\textit{Source reference.} Original lemma varieties-situation-family-divisors: sources/stacks-more-morphisms/pinned-original-source.tex, source lines 9159--9159. Cited by the retained statement or proof but not part of the collected spans; supplied only as a navigation pointer. The argument does not depend on its content.\par\medskip
\begin{lemma}
\label{lemma-bertini-irreducible}
\begin{reference}
\cite[Theorem 6.3 part 4)]{Jou}
\end{reference}
In Varieties, Situation \ref{varieties-situation-family-divisors} assume
\begin{enumerate}
\item $X$ is of finite type over $k$,
\item $X$ is geometrically irreducible over $k$,
\item there exist $v_1, v_2, v_3 \in V$ and an irreducible component
$Z$ of $H_{v_2} \cap H_{v_3}$ such that $Z \not \subset H_{v_1}$ and
$\text{codim}(Z, X) = 2$, and
\item every irreducible component $Y$ of $\bigcap_{v \in V} H_v$
has $\text{codim}(Y, X) \geq 2$.
\end{enumerate}
Then for general $v \in V \otimes_k k'$
the scheme $H_v$ is geometrically irreducible over $k'$.
\end{lemma}

\begin{proof}
In order for assumption (3) to hold, the elements $v_1, v_2, v_3$
must be $k$-linearly independent in $V$ (small detail omitted).
Thus we may choose a basis $v_1, \ldots, v_r$ of $V$ incorporating
these elements as the first $3$. Recall that
$H_{univ} \subset \mathbf{A}^r_k \times_k X$ is the ``universal divisor''.
Consider the projection $q : H_{univ} \to \mathbf{A}^r_k$
whose scheme theoretic fibres are the divisors $H_v$. By
Lemma \ref{lemma-geom-irreducible-generic-fibre}
it suffices to show that the generic fibre of $q$ is geometrically
irreducible.
To prove this we may replace $X$ by its reduction, hence
we may assume $X$ is an integral scheme of finite type over $k$.

\medskip\noindent
Let $U \subset X$ be a nonempty affine open such that
$\mathcal{L}|_U \cong \mathcal{O}_U$. Write $U = \Spec(A)$.
Denote $f_i \in A$ the element corresponding to section $\psi(v_i)|_U$
via the isomorphism $\mathcal{L}|_U \cong \mathcal{O}_U$.
Then $H_{univ} \cap (\mathbf{A}^r_k \times_k U)$ is given by
$$
H_U = \Spec(A[x_1, \ldots, x_r]/(x_1f_1 + \ldots + x_rf_r))
$$
By our choice of basis we see that $f_1$ cannot be zero because this
would mean $v_1 = 0$ and hence $H_{v_1} = X$ which contradicts
assumption (3). Hence $\sum x_if_i$ is a nonzerodivisor in
$A[x_1, \ldots, x_r]$.
It follows that every irreducible component of $H_U$ has dimension
$d + r - 1$ where $d = \dim(X) = \dim(A)$. If $U' = U \cap D(f_1)$
then we see that
$$
H_{U'} =
\Spec(A_{f_1}[x_1, \ldots, x_r]/(x_1f_1 + \ldots + x_rf_r)) \cong
\Spec(A_{f_1}[x_2, \ldots x_r]) =
\mathbf{A}^{r - 1}_k \times_k U'
$$
is irreducible. On the other hand, we have
$$
H_U \setminus H_{U'} =
\Spec(A/(f_1)[x_1, \ldots, x_r]/(x_2f_2 + \ldots + x_rf_r))
$$
which has dimension at most $d + r - 2$. Namely, for $i \not = 1$
the scheme $(H_U \setminus H_{U'}) \times_U D(f_i)$ is either empty
(if $f_i = 0$) or by the same argument as above isomorphic to an $r - 1$
dimensional affine space over an open of $\Spec(A/(f_1))$ and hence has
dimension at most $d + r - 2$.
On the other hand, $(H_U \setminus H_{U'}) \times_U V(f_2, \ldots, f_r)$
is an $r$ dimensional affine space over $\Spec(A/(f_1, \ldots, f_r))$
and hence assumption (4) tells us this has dimension at most
$d + r - 2$. We conclude that $H_U$ is irreducible for every $U$ as above.
It follows that $H_{univ}$ is irreducible.

\medskip\noindent
Thus it suffices to show that the generic point of $H_{univ}$
is geometrically irreducible over the generic point of $\mathbf{A}^r_k$, see
Varieties, Lemma
\ref{varieties-lemma-geometrically-irreducible-function-field}.
Choose a nonempty affine open $U = \Spec(A)$ of $X$
contained in $X \setminus H_{v_1}$ which meets the irreducible
component $Z$ of $H_{v_2} \cap H_{v_3}$ whose existence is asserted in
assumption (3). With notation as above we have to prove that
the field extension
$$
\text{Frac}(A[x_1, \ldots, x_r]/(x_1f_1 + \ldots + x_rf_r))/
k(x_1, \ldots , x_r)
$$
is geometrically irreducible. Observe that $f_1$ is invertible in $A$
by our choice of $U$.
Set $K = \text{Frac}(A)$ equal to the fraction field of $A$.
Eliminating the variable $x_1$ as above,
we find that we have to show that the field extension
$$
K(x_2, \ldots, x_r)/
k(x_2, \ldots, x_r, -\sum\nolimits_{i = 2, \ldots, r} f_1^{-1}f_i x_i)
$$
is geometrically irreducible. By Lemma \ref{lemma-amazing-bertini-lemma}
it suffices to show that for some $c_2, \ldots, c_r \in k$ the elements
$$
x_2, \ldots, x_r, \sum\nolimits_{i = 2, \ldots, r} f_1^{-1}f_i x_i,
\sum\nolimits_{i = 2, \ldots, r} c_if_1^{-1}f_i
$$
are algebraically independent over $k$ in the fraction field of
$A[x_2, \ldots, x_r]$. This follows from
Lemma \ref{lemma-algebraically-independent}
and the fact that $Z \cap U$ is an irreducible
component of $V(f_1^{-1}f_2, f_1^{-1}f_3) \subset U$.
\end{proof}

\begin{thebibliography}{1}
\bibitem{Jou}
Jouanolou, Jean-Pierre.
\newblock {\em Th\'{e}or\`{e}mes de Bertini et Applications}.
\newblock Progress in Mathematics 42. Birkh\"{a}user, 1983.
\end{thebibliography}
\end{document}
